\chapter{Corners and the scales in graphical estimates}\label{sec:examples-graphical-cutoff-scales}
These examples connect the local estimates of \cite{EH91} with the explicit expanding curves constructed in Section \ref{sec:examples-expanders}. They distinguish a sharp scaling exponent from a claim that every numerical constant in an estimate is optimal.

\begin{example}[The time exponent one half cannot be improved for Lipschitz graphs]\label{ex:graphical-half-time-sharpness}
For the expanding curve $U(x,t)=\sqrt t\,u_a(x/\sqrt t)$ of Proposition \ref{prop:expander-corner-initial-trace}, the initial trace is $\ell_a|x|$ and
\[
 U(0,t)-U(0,0)=a\sqrt t.
\]
Consequently a uniform time estimate with exponent $\alpha>1/2$ at zero cannot hold for all Lipschitz initial graphs, even for a single such nontrivial flow with its own finite constant. The curvature satisfies $|A|(0,t)=a/(2\sqrt t)$, so the positive-time curvature scale also occurs in an exact solution.
\end{example}
\begin{proof}
The profile has $u_a(0)=a>0$, giving the height identity. Dividing it by $t^\alpha$ gives $a t^{1/2-\alpha}\to\infty$, excluding a finite constant for an estimate of that exponent near zero. Proposition \ref{prop:expander-curvature-example} gives the curvature identity. These examples have bounded slope at all positive times by Lemma \ref{lem:expander-global-slope}; the singular initial corner is why an initial curvature bound is unavailable. Products with Euclidean factors give the same height and curvature exponents in every hypersurface dimension.
\end{proof}

\begin{lemma}[The tip height parametrizes all positive symmetric cone slopes]\label{lem:expander-tip-slope-bijection}
For the even convex expander profiles, the map $a\mapsto\ell_a$ is a continuous strictly increasing bijection from $(0,\infty)$ to $(0,\infty)$. Moreover
\[
 \frac{a}{\sqrt2}\leq\ell_a,
 \qquad \arctan\ell_a\leq\frac{\sqrt\pi}{2}a.
\]
\end{lemma}
\begin{proof}
The time estimate in Theorem \ref{thm:graph-lipschitz-half-holder-time}, applied at the tip to the explicit expander with initial Lipschitz constant $\ell_a$ and $n=1$, gives $a\leq\sqrt2\ell_a$. For the upper angle estimate, the defect equation of Proposition \ref{prop:expander-graph-defect} gives
\[
 (\arctan p)'=q/2\leq(a/2)e^{-x^2/4}.
\]
Integrate from zero to infinity and use $p(0)=0$, $p(\infty)=\ell_a$ and $\int_0^\infty e^{-x^2/4}\,dx=\sqrt\pi$.

If $a<b$ but $\ell_a\geq\ell_b$, their expanding flows have ordered initial traces in the opposite order. Theorem \ref{thm:entire-bounded-slope-comparison} gives $U_b\leq U_a$, hence $b\leq a$ at the tip at time one, a contradiction. Thus the slope is strictly increasing.

To check continuity rather than assuming it from monotonicity, fix a compact interval of positive parameters. Smooth ODE dependence makes $p_a(1)$ and $q_a(1)$ bounded and continuous there. Lemma \ref{lem:expander-global-slope} bounds $p_a(x)$ uniformly for $x\geq1$. The defect bound therefore gives $p_a'(x)\leq C e^{-x^2/4}$ uniformly over those parameters. Hence $p_a(R)$ converges uniformly to $\ell_a$ as $R\to\infty$. Every finite-$R$ value is continuous by ODE dependence, so its uniform limit is continuous. The angle bound implies $\ell_a\to0$ as $a\downarrow0$, and the lower bound implies $\ell_a\to\infty$ as $a\to\infty$. Continuity and strict increase finish the bijection proof.
\end{proof}

\begin{theorem}[The unique bounded-slope graphical expander for a symmetric corner]\label{thm:prescribed-symmetric-corner-expander}
For every $\ell>0$ there is a unique smooth entire graphical self-expander whose slope is bounded and whose dilation flow converges locally uniformly to $\ell|x|$ at time zero. It is the even strictly convex profile $u_a$ with $\ell_a=\ell$. Uniqueness here is in this graphical bounded-slope class, not among arbitrary immersed networks or weak flows. The planar expander setting is studied in \cite{Hal12}, Section 6; the uniqueness proof here uses graphical comparison.
\end{theorem}
\begin{proof}
Lemma \ref{lem:expander-tip-slope-bijection} supplies one parameter $a$. Theorem \ref{thm:convex-graphical-expander} and Proposition \ref{prop:expander-corner-initial-trace} supply the profile and its uniform initial trace. Any other profile in the stated class has a dilation flow smooth for positive time with the same locally uniform initial trace; dilation preserves the bound on its graph slope. Theorem \ref{thm:entire-bounded-slope-comparison} gives equality of the two flows. Their time-one slices are the profiles themselves. Evenness and strict convexity are therefore conclusions for every profile in this class, rather than extra assumptions imposed to prove uniqueness.
\end{proof}

\begin{example}[The spatial-radius term in local curvature estimates is necessary]\label{ex:scherk-local-curvature-scale}
For $\lambda>0$, consider the rescaled Scherk graph
\[
 u_\lambda(x,y)=\lambda\bigl(\log\cos(y/\lambda)-\log\cos(x/\lambda)\bigr).
\]
It is a stationary graphical flow on $|x|,|y|<\pi\lambda/2$. On the base ball $B_\lambda(0)$ its slope is bounded independently of $\lambda$, whereas $|A|^2(0)=2/\lambda^2$ at every time. Thus the $R^{-2}$ term in Theorem \ref{thm:graphical-base-interior-curvature} cannot be omitted when $R$ is finite, and its power agrees with scaling.
\end{example}
\begin{proof}
Example \ref{ex:scherk-graph} verifies the minimal graph equation at scale one. Euclidean dilation preserves minimality, hence stationarity under flow. On $B_\lambda$, both dimensionless coordinates have magnitude at most one, so $|Du_\lambda|^2\leq2\tan^2(1)$. At the origin the gradient is zero and the Hessian is $\operatorname{diag}(1/\lambda,-1/\lambda)$. The graph metric is Euclidean there, giving $|A|^2=2/\lambda^2$. If the local curvature bound involved only $C(V)/t$, sending $t\to\infty$ along this stationary flow would force its curvature to vanish, a contradiction. The ball lies strictly inside the graph domain, so failure is not due to a lack of smoothness on its closure.
\end{proof}
